\section{Certificate Architecture and Independent Verification}
\label{sec:verification}

Large symbolic calculations are difficult to trust by inspection.  We use a
proof-carrying architecture: discovery may be complex, but the final claim is
accepted only after a narrower verifier rechecks explicit evidence.

\subsection{Trusted and untrusted components}

\begin{figure}[t]
\centering
\resizebox{\textwidth}{!}{%
\begin{tikzpicture}[
 node distance=6mm and 10mm,
 box/.style={draw,rounded corners,align=center,minimum height=9mm,minimum width=33mm},
 untrusted/.style={box,fill=orange!10},
 trusted/.style={box,fill=green!10},
 arr/.style={-{Stealth[length=2mm]},thick}
]
\node[untrusted] (gen) {symbolic generator\\and grouping search};
\node[untrusted,below=of gen] (hpc) {parallel/HPC discovery\\and cached ledgers};
\node[box,right=of gen,yshift=-8mm,fill=blue!5] (bundle) {hashed certificate bundle\\JSON, groups, summaries};
\node[trusted,right=of bundle] (fast) {fast verifier\\schema + arithmetic};
\node[trusted,below=of fast] (deep) {deep replay\\algebra + group membership};
\node[trusted,below=of deep] (ref) {reference verifier\\standard library only};
\node[trusted,right=of deep] (claim) {accepted claim\\or explicit failure};
\draw[arr] (gen)--(bundle);
\draw[arr] (hpc)--(bundle);
\draw[arr] (bundle.east)--(fast.west);
\draw[arr] (bundle.east)--++(2mm,0)|-(deep.west);
\draw[arr] (bundle.east)--++(2mm,0)|-(ref.west);
\draw[arr] (fast.east)--++(5mm,0)|-(claim.north);
\draw[arr] (deep.east)--(claim.west);
\draw[arr] (ref.east)--++(5mm,0)|-(claim.south);
\end{tikzpicture}
}
\caption{Trust boundary.  Parallel discovery and heuristic grouping are
untrusted producers.  The accepted result depends on explicit evidence and
three verification paths with different trust surfaces.}
\label{fig:trust-boundary}
\end{figure}

The trusted claim comprises the mathematical lemmas in this paper, the small
verification programs, exact integer/rational arithmetic, and the cryptographic
binding between manifest and sidecars.  It does not require trusting the job
scheduler, the heuristic's optimality, cached stdout, or a claimed successful
HPC run.

\subsection{Certificate contents}

The schema-v3 certificate records:
\begin{itemize}
  \item the complete benchmark tuple, normalization, formula, time, tolerance,
        and primary resource metric;
  \item the pinned theorem instantiation used for the baseline;
  \item coefficient precision policies and rational interval endpoints;
  \item D4 and D5 group-certificate paths, SHA-256 digests, term counts, group
        counts, maximum sizes, and norm enclosures;
  \item D4--D7 and tail contributions at the accepted integer;
  \item the accepted and previous-step global-error rationals;
  \item integer resource arithmetic and the exact improvement ratio; and
  \item auxiliary audit results explicitly separated from the primary claim.
\end{itemize}
Large membership data are sidecars rather than embedded in prose.  The
manifest hash prevents a reviewer from inadvertently validating a different
file with the same name.

\subsection{Fast and deep modes}

Fast verification checks the schema, paths, hashes, integer identities,
rational comparisons, resource formulas, and headline invariants.  It also
reconstructs the formula constants and the complete D4--D7 plus tail ledger at
both the accepted and adjacent step, rather than trusting a self-consistent
submitted sum.  It verifies sidecar coverage, pairwise anticommutation, and
rational square-root enclosures.  Deep verification additionally regenerates
all 31 baseline centers and every D4/D5 coefficient from the Suzuki formula,
then requires exact interval equality with the sidecars.

The two modes serve different purposes.  Fast verification answers whether a
checked-in bundle contains valid sidecar norm witnesses and a fully reproduced
finite-step ledger.  Deep verification answers whether the large coefficient
maps themselves arise from the declared Suzuki formula.  Neither mode treats a
rounded decimal in a PDF as evidence.

A third, deliberately smaller program uses only the Python standard library
and imports no \code{trottercert} module.  It independently implements
canonical sidecar hashing, D4/D5 coverage, same-support and symplectic
anticommutation checks, rational square-root domination, the Suzuki tail,
finite-step integration, the 95/94 comparison, and resource identities.  It
does not regenerate the large coefficient maps; that distinct upstream
obligation remains with deep mode.  Agreement therefore cross-checks the
downstream proof logic without falsely claiming fully independent symbolic
generation.

\subsection{Negative tests and failure localization}

A certificate system is useful only if nearby corruptions fail.  The test
suite mutates representative fields: a group member, a coefficient endpoint,
the accepted step count, D6, D7, the tail, the adjacent-step value, a resource
total, and a sidecar digest.  Ledger attacks also recompute the submitted total
to ensure that mere internal consistency is insufficient.  Each mutation must
cause a nonzero verifier result.  Further unit tests cover Pauli phase
multiplication, symplectic commutation, periodic canonicalization, rational
square-root enclosure, and the local 16-Pauli lemma.

Failures are reported by layer.  A hash mismatch is distinguished from a
schema error; a duplicate group member from a non-anticommuting pair; and a
regenerated tail mismatch from an adjacent-step comparison.  This makes independent
reproduction tractable: a reviewer need not rerun the entire discovery
pipeline to diagnose a damaged artifact.

\subsection{Claim classes}

We label statements by evidence class:
\begin{description}
  \item[Certified.] Derived from a manifest-bound witness and accepted by the
  verifier, such as the 95-step result and integer resource counts.
  \item[Cross-checked.] Supported by a structurally independent small-system or
  exhaustive local calculation, but not the main large-system proof.
  \item[Conditional.] An arithmetic consequence under a stated future premise,
  such as the norm budget required for a fivefold result.
  \item[Prospective.] A research direction or engineering optimization with no
  claimed current performance.
\end{description}
This vocabulary prevents an observed trend, a possible optimization, and a
completed theorem from being conflated in the abstract or conclusion.
